\section{The Employment Counterfactual}

\subsection{The Opportunity Cost of Foregone Earnings}

A founder who says, ``If I work for myself, I keep the upside,'' is making only half a comparison. The other half is: What would I have earned, learned, saved, invested, and been able to pursue if I had remained employed?

Consider a hypothetical engineer with a career path that allows them to save and invest ₹12 lakh a year. Suppose the savings earn an illustrative annual return of 8\%, with each year's contribution made at year-end. After 10 years, the account would contain approximately:
\begin{equation}
  S_{10} = ₹12\text{ lakh} \times \frac{(1 + 0.08)^{10} - 1}{0.08} \approx ₹1.74\text{ crore} \label{eq:counterfactual-savings}
\end{equation}

This is an arithmetic illustration, not a promise of investment returns. It excludes taxes, fees, inflation, and changes in earnings. Some professionals could save much less; others substantially more. Still, the example exposes a real opportunity cost: a good job can be a powerful wealth-building engine when paired with disciplined saving and ownership of financial assets.

Compare this with the hypothetical venture payoff introduced earlier: 70\% chance of zero equity proceeds, 25\% chance of ₹3 crore, and 5\% chance of ₹30 crore, producing an expected gross payoff of ₹2.25 crore. One might conclude that the startup is better because ₹2.25 crore exceeds ₹1.74 crore. That conclusion would be premature.

The founder may receive salary or distributions during the 10 years, or may receive little. They may contribute personal capital, borrow money, suffer years of weak income, or remain employed for part of the period. Proceeds may be taxed differently from employment income. The startup payout is highly variable and potentially illiquid, while the investment portfolio follows its own uncertain return path. And the probabilities of the venture outcomes were invented solely to illustrate distributional reasoning. They were not measured from a relevant cohort.

A proper comparison would model both pathways as distributions of after-tax household wealth over the same period, include earnings and savings during each path, account for different time demands and risks, and consider the ability to switch paths. For many people this calculation would favor employment. For some it would favor a business. For others, especially those with a strong employment position and modest initial business costs, the best choice might be a hybrid sequence.

\subsection{Leverage Within Established Institutions}

There is a further reason not to dismiss employment: large organizations can provide substantial leverage. A senior researcher may direct a scientific team, gain access to expensive infrastructure, and influence products used around the world. A highly compensated executive may earn equity incentives and be responsible for a business unit larger than most startups. A professional with scarce expertise may achieve both financial independence and meaningful impact without assuming founder risk.

\textbf{A company is not the only way to gain leverage. It is one way to choose and own the leverage you build.}

This is why the most useful question is not ``Can I become wealthy as an employee?'' Of course some people can. Nor is it ``Can I become a billionaire as a founder?'' In rare cases, yes. The question is:

\textbf{Given my abilities, current position, financial obligations, appetite for uncertainty, and ambitions, which path (or sequence of paths) offers the most attractive combination of expected outcomes and the life I actually want to live?}

One person's answer may be to found immediately. Another's may be to save for 5 years, become exceptionally good at a domain, and then launch. Another's may be to build a business during evenings and weekends, with clear boundaries and no conflict with employment duties. Another's may be to take an entrepreneurial role inside an established firm. Being serious about entrepreneurship does not require treating every salaried year as a failure.
